%%%%%%%%%%%%%%%%%%%%%%%%%%%%%%%%%%%%%%%%%%%%%%%%%%%%%%%%%%%%%%%%%%%%%%
% Lớp 10 | Xác suất | Xác suất của biến cố
% Dạng 0D0?2-9: Liên quan vấn đề khác
% Số câu: 18  (sinh tự động bởi tach_cau.py; từ nay thầy sửa trực tiếp file này)
%%%%%%%%%%%%%%%%%%%%%%%%%%%%%%%%%%%%%%%%%%%%%%%%%%%%%%%%%%%%%%%%%%%%%%

%%%%%%%%%%%%%%%% TRẮC NGHIỆM 4 PHƯƠNG ÁN %%%%%%%%%%%%%%%%

% nguồn: MaHoaNH_TracNghiem_Lop10
\begin{ex}%[0D0N2-9]%[ma:66b14613]
	Trong không gian cho $10$ điểm, trong đó có $A$ và $B$,  sao cho không có ba điểm nào thẳng hàng. Quỳnh vẽ một đoạn thẳng nối hai điểm tùy ý trong $10$ điểm. Tính xác suất để đoạn thẳng mà Quỳnh vẽ là đoạn $AB$.
	\choice 
		 {$\dfrac{1}{10}$}
		 {$\dfrac{1}{90}$}
		 {$\dfrac{1}{5}$}
		{\True  $\dfrac{1}{45}$}
	 \loigiai{
		Tổng số đoạn thẳng có thể vẽ là $\mathrm{C}_{10}^2 =  45$.\\
		Chỉ có duy nhất $1$ đoạn là $AB$ $\Rightarrow \mathrm{P}(A) =  \dfrac{1}{45}$.
	 }
\end{ex}

% nguồn: MaHoaNH_TracNghiem_Lop10
\begin{ex}%[0D0N2-9]%[ma:b0974785]
	Chọn ngẫu nhiên một quân bài trong bộ bài tú lơ khơ gồm $52$ lá. Tính xác suất để quân bài  được chọn có số nút là $9$.
	\choice 
		{\True  $\dfrac{1}{13}$}
		 {$\dfrac{1}{12}$}
		 {$\dfrac{1}{11}$}
		{$\dfrac{1}{10}$} 
	 \loigiai{
		Số phần tử của không gian mẫu là $n(\Omega)  = \mathrm{C}_{52}^1 = 52$.\\
		Gọi $A$ là biến cố \lq\lq Quân bài được chọn có  số nút là $9$\rq\rq.\\
		Mỗi  giá trị nút có $4$ quân bài $\Rightarrow n(A) = 4\Rightarrow \mathrm{P}(A) = \dfrac{4}{52} = \dfrac{1}{13}$.
	} 
\end{ex}

% nguồn: MaHoaNH_TracNghiem_Lop10
\begin{ex}%[0D0N2-9]%[ma:f00f7063]
	Từ một đội văn nghệ gồm $5$ nam và $8$ nữ cần lập một  nhóm gồm $4$ người hát tốp ca. Xác suất để trong $4$ người được chọn đều là nam bằng
	 \choice
		{$\dfrac{\mathrm{C}_8^4}{\mathrm{C}_{13}^4}$} 
		 {$\dfrac{\mathrm{A}_5^4}{\mathrm{C}_8^4}$}
		{\True $\dfrac{\mathrm{C}_5^4}{\mathrm{C}_{13}^4}$} 
		{$\dfrac{\mathrm{C}_8^4}{\mathrm{A}_{13}^4}$} 
	 \loigiai{
		Tổng số cách chọn $4$ người từ $13$ người $n(\Omega) =  \mathrm{C}_{13}^4$.\\
		Gọi $A$ là biến cố \lq\lq Chọn  được $4$ người đều là nam\rq\rq.\\
		Vì có $5$ nam $\Rightarrow n(A)  = \mathrm{C}_5^4$.\\
		Vậy xác suất của biến cố  $A$ là $\mathrm{P}(A) = \dfrac{\mathrm{C}_5^4}{\mathrm{C}_{13}^4}$.
	 }
\end{ex}

% nguồn: Ok_TracNghiemLop11_Lan1
\begin{ex}%[0D0H2-9]%[ma:6d61bfce]
	Hai xạ thủ cùng bắn vào bia, mỗi xạ thủ bắn một lần. Xác suất bắn trúng đích của xạ thủ thứ nhất là $0{,}8$ và xạ thủ thứ hai là $0{,}9$. Tính xác suất  của biến cố \lq\lq Có đúng một xạ thủ bắn trúng đích\rq\rq.
	 \choice
		{$0{,}18$} 
		{\True $0{,}26$} 
		{$0{,}08$} 
		{$0{,}72$} 
	\loigiai{ 
	Gọi $A$ là biến cố \lq\lq Xạ thủ thứ nhất bắn trúng đích\rq\rq, suy ra $\mathrm{P}(A)=0{,}8$. \\ 
	Khi đó $\overline{A}$ là biến cố \lq\lq Xạ thủ thứ nhất bắn trượt\rq\rq,  suy ra $\mathrm{P}\left(\overline{A}\right)=0{,}2$. \\
	Gọi $B$ là biến cố \lq\lq Xạ thủ  thứ hai bắn trúng đích \rq\rq, suy ra $\mathrm{P}(B)=0{,}9$. \\
	Khi đó $\overline{B}$  là biến cố \lq\lq Xạ thủ thứ hai bắn trượt \rq\rq, suy ra $\mathrm{P}\left(\overline{B}\right)=0{,}1$. \\
	Gọi  $C$ là biến cố \lq\lq Có đúng một xạ thủ bắn trúng đích\rq\rq.\\
	Suy ra $C=A\overline{B}\cup\overline{A}B$. \\ 
	Vì $A$ và $B$ là hai biến cố độc lập  nên $\mathrm{P}(C)=\mathrm{P}\left(A\overline{B}\right)+\mathrm{P}\left(\overline{A}B\right)=0{,}26$. \\
	Vậy xác suất của biến cố \lq\lq Có  đúng một xạ thủ bắn trúng đích\rq\rq\ là $0{,}26$.
	 }
\end{ex}

% nguồn: MaHoaNH_TracNghiem_Lop10
\begin{ex}%[0D0V2-9]%[ma:300a9464]
	Gieo một con súc sắc đồng  chất $2$ lần. Tính xác suất để tổng $2$ mặt thu được của $2$ lần gieo là số lẻ.
	\choice 
		 {$\dfrac{12}{36}$}
		 {$\dfrac{25}{36}$}
		 {$\dfrac{6}{36}$}
		 {\True $\dfrac{18}{36}$}
	 \loigiai{
		Ta có $n(\Omega) = 6 \cdot 6 =  36$.\\
		Gọi $A$ là  biến cố \lq\lq Tổng $2$ mặt của $2$ lần gieo là lẻ\rq\rq.\\
		Vì  mỗi số chẵn ở lần tung đầu tiên tương ứng với một số lẻ ở lần tung thứ hai và ngược lại thì ta được tổng $2$ lần tung là $1$ số lẻ.\\
		Như vậy với $3$ số chẵn lần đầu ứng với $3$ số lẻ lần sau $\Rightarrow$ có $3 \cdot 3 = 9$  cách gieo.\\
		Với $3$ số lẻ lần đầu ứng với $3$ số chẵn lần sau cũng có $3 \cdot 3 = 9$ cách gieo  $\Rightarrow \mathrm{P}(A) = \dfrac{18}{36} = 0{,}5$.
	 }
\end{ex}

% nguồn: MaHoaNH_TracNghiem_Lop10
\begin{ex}%[0D0V2-9]%[ma:b6b2678a]
	Bốn bạn nam và bốn bạn nữ được xếp ngồi ngẫu nhiên vào $8$ ghế xếp thành hai dãy đối diện nhau. Xác suất sao cho nữ ngồi đối diện nhau  là
	\choice 
		 {$\dfrac{3}{4}$}
		 {$\dfrac{1}{2}$}
		{$\dfrac{5}{6}$} 
		 {\True $\dfrac{3}{35}$}
	 \loigiai{
		Xếp một bạn nữ vào $1$ trong $8$ ghế có $8$  cách.\\
		Tiếp theo chọn $1$ trong $3$ bạn nữ còn lại để xếp  vào ghế đối diện có $3$ cách.\\
		Tiếp tục xếp bạn nữ thứ $3$ vào $1$ trong $6$ ghế còn lại có $6$ cách, bạn nữ còn lại buộc phải ngồi ghế đối diện có  $1$ cách.\\
		$4$ ghế trống còn lại xếp $4$ bạn nam vào có $4!$ cách.\\ 
		Gọi $A$ là biến cố \lq\lq Nữ ngồi  đối diện nhau\rq\rq.\\
		Số trường hợp  thuận lợi là $n(A) = 8 \cdot 3 \cdot 6 \cdot 1 \cdot 4! = 3456$.\\
		Số phần tử của không gian mẫu là $n(\Omega) = 8!\Rightarrow \mathrm{P}(A) =  \dfrac{3456}{8!} = \dfrac{3}{35}$.
	} 
\end{ex}

% nguồn: MaHoaNH_TracNghiem_Lop10
\begin{ex}%[0D0V2-9]%[ma:2a939351]
	Một hộp đựng $11$ tấm thẻ được đánh số từ $1$ đến $11$. Chọn ngẫu nhiên $6$ tấm thẻ. Gọi $P$ là xác suất  để tổng số ghi trên $6$ tấm thẻ ấy là một số lẻ. Khi đó $\mathrm{P}(A)$ bằng
	\choice 
		 {$\dfrac{100}{231}$}
		 {$\dfrac{115}{231}$}
		{$\dfrac{1}{2}$} 
		{\True $\dfrac{118}{231}$} 
	\loigiai{ 
		Số phần tử của không gian mẫu  là $n(\Omega)=\mathrm{C}_{11}^6 = 462$.\\
		Gọi $A$ là biến cố: \lq\lq Tổng số ghi  trên $6$ tấm thẻ là một số lẻ\rq\rq.\\
		Từ $1$ đến $11$  có $6$ số lẻ và $5$ số chẵn. Để có tổng là một số lẻ ta có ba trường hợp:\\
		\textbf{Trường hợp 1:} Chọn $1$ thẻ mang  số lẻ và $5$ thẻ mang số chẵn có $\mathrm{C}_6^1\cdot \mathrm{C}_5^5 = 6$ cách.\\
		\textbf{Trường hợp 2:} Chọn $3$ thẻ lẻ và $3$ thẻ chẵn có $\mathrm{C}_6^3\cdot \mathrm{C}_5^3 = 200$  cách.\\
		\textbf{Trường hợp 3:} Chọn $5$ thẻ lẻ và $1$ thẻ  chẵn có $\mathrm{C}_6^5\cdot \mathrm{C}_5^1 = 30$ cách.\\
		 Do đó $n(A)=6+200+30=236$ nên $\mathrm{P}(A)=\dfrac{236}{462}=\dfrac{118}{231}$.
	 }
\end{ex}

% nguồn: MaHoaNH_TracNghiem_Lop10
\begin{ex}%[0D0V2-9]%[ma:645add6b]
	 Một hộp có $5$ viên bi xanh, $6$ viên bi đỏ và $7$ viên bi vàng. Chọn ngẫu nhiên $5$ viên bi trong hộp. Xác suất để $5$ viên bi được chọn chỉ có một màu là
	 \choice
		{\True  $\dfrac{1}{306}$}
		 {$\dfrac{1}{1428}$}
		{$\dfrac{1}{408}$} 
		{$\dfrac{1}{8568}$} 
	 \loigiai{
		Số phần tử của không gian mẫu là $n(\Omega) = \mathrm{C}_{18}^5 =  8568$.\\
		Gọi $A$ là biến cố \lq\lq  $5$ viên bi được chọn chỉ có một màu\rq\rq.
		\begin{itemize}
			\item Chọn $5$ bi đỏ: $\mathrm{C}_6^5 = 6$  cách.
			\item Chọn $5$ bi vàng: $\mathrm{C}_7^5 = 21$ cách. 
			\item Chọn $5$ bi xanh: chỉ có $1$ cách (vì chỉ có đúng $5$  viên).
		\end{itemize}
		$\Rightarrow n(A) = 6 + 21 + 1 = 28 \Rightarrow \mathrm{P}(A)  = \dfrac{28}{8568} = \dfrac{1}{306}$.
	 }
\end{ex}

% nguồn: MaHoaNH_TracNghiem_Lop10
\begin{ex}%[0D0V2-9]%[ma:6d8f7d1c]
	Gieo $2$ con xúc xắc cân đối và đồng chất. Xác suất để tổng số chấm xuất hiện trên hai mặt  của $2$ con xúc xắc đó không vượt quá $4$ là
	 \choice
		{$\dfrac{5}{18}$} 
		{$\dfrac{2}{9}$} 
		{$\dfrac{2}{3}$} 
		 {\True $\dfrac{1}{6}$}
	 \loigiai{
		Số  phần tử của không gian mẫu là $n(\Omega) = 6^2 = 36$.\\
		Gọi $A$ là biến cố \lq\lq  Tổng số chấm của hai xúc xắc không quá $4$\rq\rq.\\
		Các trường hợp thỏa mãn  là
		 \[
		(1;1), (1;2), (1;3), (2;1), (2;2),  (3;1)
		\Rightarrow n(A) =  6.
		 \]
		Vậy $\mathrm{P}(A)  = \dfrac{n(A)}{n(\Omega)} = \dfrac{6}{36} = \dfrac{1}{6}$.
	} 
\end{ex}

% nguồn: MaHoaNH_TracNghiem_Lop10
\begin{ex}%[0D0V2-9]%[ma:a8b40727]
	Một nhóm gồm $8$ nam và $7$ nữ. Chọn ngẫu  nhiên $5$ bạn. Xác suất để trong $5$ bạn được chọn có cả nam lẫn nữ mà nam nhiều hơn nữ là
	 \choice
		 {$\dfrac{60}{143}$}
		{\True  $\dfrac{238}{429}$}
		{$\dfrac{210}{429}$} 
		 {$\dfrac{82}{143}$}
	\loigiai{ 
		Số phần tử của không gian mẫu là $n(\Omega) =  \mathrm{C}_{15}^5 = 3\,003$.\\
		Xét các trường  hợp trong đó số nam nhiều hơn nữ:
		\begin{itemize}
			\item Trường hợp 1: $4$ nam và $1$ nữ: $\mathrm{C}_8^4 \cdot \mathrm{C}_7^1 = 70  \cdot 7 = 490$.
			\item Trường hợp 2: $3$ nam và  $2$ nữ: $\mathrm{C}_8^3 \cdot \mathrm{C}_7^2 = 56 \cdot 21 = 1\,176$.
		\end{itemize}
		$\Rightarrow n(A) = 490 + 1\,176 = 1\,666 
		 \Rightarrow \mathrm{P}(A) =\dfrac{1\,666}{3\,003} = \dfrac{238}{429}$.
	} 
\end{ex}

% nguồn: MaHoaNH_TracNghiem_Lop10
\begin{ex}%[0D0V2-9]%[ma:502facf7]
	Cho tập $A = \{1; 2; 3; 4; 5; 6\}$.  Từ tập $A$ có thể lập được bao nhiêu số tự nhiên có $3$ chữ số khác nhau. Tính xác suất biến cố sao cho tổng $3$ chữ số bằng $9$.
	 \choice
		{$\dfrac{1}{20}$} 
		{\True  $\dfrac{3}{20}$}
		 {$\dfrac{9}{20}$}
		 {$\dfrac{7}{20}$}
	\loigiai{ 
		Gọi $A$ là  biến cố \lq\lq Số tự nhiên có tổng $3$ chữ số bằng $9$\rq\rq.\\
		Số số tự nhiên có $3$ chữ số khác nhau được lập  từ $A$ là $A_6^3 = 120$.\\
		Số phần tử của không gian mẫu $n(\Omega) = 120$.\\ 
		Tổng $3$ chữ số bằng $9$ có thể là 
		\[ 
		1+2+6;\quad  1+3+5;\quad 2+3+4.
		\] 
		Mỗi tổ hợp có  $3! = 6$ hoán vị $\Rightarrow n(A) = 18\Rightarrow \mathrm{P}(A) = \dfrac{18}{120} = \dfrac{3}{20}$.
	} 
\end{ex}

% nguồn: MaHoaNH_TracNghiem_Lop10
\begin{ex}%[0D0V2-9]%[ma:1b84a50b]
	Trò chơi bầu cua tôm cá là một trò chơi khá thông dụng. Vì đây là trò đỏ đen nên thường được chơi trong những ngày Tết để mọi người thử vận may. Luật chơi như sau: Có $3$ quân súc sắc trên đó thay vì ghi các số $1,2,3,4,5,6$ thì có các hình Bầu, Cua, Tôm, Cá, Nai, Gà. Bàn chơi là một mảnh bìa trên đó cũng ghi hình $6$ con vật nói trên. Ba con hàng phía trên là Nai, Bầu, Gà (trên cạn), ba con hàng phía dưới là Tôm, Cua, Cá (dưới nước). Mỗi một lần chơi, người chơi chọn một hoặc nhiều ô để đặt tiền vào đó. Nhà cái sẽ dùng một cái đĩa, để $3$ quân súc sắc trên đó, úp lại bằng một cái bát và xóc đi xóc lại nhiều lần, sau đó mở bát ra. Các mặt ngửa lên của các quân súc sắc chính là các mặt thắng. Giả sử bạn đặt $x$ đồng vào cửa Cua mà có $1$ con cua xuất hiện thì bạn sẽ được trả thêm $x$ đồng; có $2$ con Cua xuất hiện thì bạn được trả thêm $2x$ đồng; có $3$ Con cua xuất hiện thì bạn được trả thêm  $3x$ đồng. Nếu không có con Cua nào xuất hiện thì bạn thua cuộc và bị mất $x$ đồng ấy. Xác suất để Bạn Nam đặt $20\,000$ đồng vào cửa Nai trong $1$ lần chơi và thắng được $40\,000$ đồng, sau đó bạn $A$ tiếp tục đặt $40\,000$ đồng tiếp tục vào cửa Nai thắng được $120\,000$ đồng là bao nhiêu?
	 \choice
		{$\dfrac{2}{27}$} 
		 {$\dfrac{1}{216}$}
		 {$\dfrac{5}{72}$}
		{\True $\dfrac{5}{15\,552}$} 
	\loigiai{ 
		Gọi $A$ là biến  cố bạn Nam đặt $20\,000$ đồng và thắng được $40\,000$ đồng.\\
		Khi đó, trong $3$ con súc sắc phải xuất hiện đúng $2$  mặt Nai.\\
		Do đó ta được xác suất thắng cuộc của bạn Nam trong trường hợp này là $\mathrm{P} 
		(A)=3\cdot \dfrac{1}{6}\cdot  \dfrac{1}{6}\cdot \dfrac{5}{6}=\dfrac{5}{72}$.\\
		 Gọi $B$ là biến cố bạn Nam đặt $40\,000$ đồng và thắng được $120\,000$ đồng.\\
		Khi đó, trong $3$ con súc sắc phải xuất hiện  đều mặt Nai.\\
		Do đó ta được xác suất thắng cuộc của bạn Nam  trong trường hợp này là $\mathrm{P}
		 (B)=\dfrac{1}{6}\cdot \dfrac{1}{6}\cdot \dfrac{1}{6}=\dfrac{1}{216}$.\\
		Vậy xác suất  để bạn Nam chiến thắng cả $2$ lần chơi theo bài toán đã cho là\\
		$\mathrm{P}(AB)=  \mathrm{P}(A)\cdot \mathrm{P}(B)=\dfrac{5}{72}\cdot \dfrac{1}{216}=\dfrac{5}{15\,552}$.
	 }
\end{ex}

% nguồn: MaHoaNH_TracNghiem_Lop10
\begin{ex}%[0D0V2-9]%[ma:95e9ed24]
	Một người bỏ ngẫu nhiên ba lá thư vào ba chiếc phong bì đã ghi địa chỉ. Xác suất để có ít nhất một lá thư  được bỏ đúng phong bì là
	 \choice
		{$\dfrac{1}{2}$} 
		{\True  $\dfrac{2}{3}$}
		{$\dfrac{1}{3}$} 
		{$\dfrac{5}{6}$} 
	 \loigiai{
		Số phần tử của không  gian mẫu là $n(\Omega) = 3! = 6$.\\
		Gọi $A$ là biến  cố \lq\lq Có ít nhất một lá thư được bỏ đúng phong bì\rq\rq.\\
		Ta  xét các trường hợp:
		\begin{itemize}
			\item Nếu lá thư thứ nhất bỏ đúng phong bì, hai lá còn lại để sai thì có duy nhất $1$ cách. 
			\item Nếu lá thư thứ hai bỏ đúng phong bì, hai lá còn lại để sai thì cũng chỉ  có $1$ cách.
			\item Nếu lá thư thứ ba bỏ đúng phong bì, hai lá còn  lại để sai thì cũng có $1$ cách.
			\item Không thể có trường hợp hai lá thư bỏ đúng và  một lá thư sai.
			\item Cả ba lá thư đều  đúng có duy nhất $1$ cách.
		\end{itemize}
		 Vậy $n(A) = 1+1+1+1 = 4$ cách chọn thỏa mãn.\\
		$\Rightarrow  \mathrm{P}(A) = \dfrac{4}{6} = \dfrac{2}{3}$.
	 }
\end{ex}

% nguồn: MaHoaNH_TracNghiem_Lop10
\begin{ex}%[0D0V2-9]%[ma:c98521c4]
	Gọi $E$  là tập hợp các số tự nhiên gồm ba chữ số khác nhau từng đôi một được chọn từ các số $0,1,2,3,4,5,6,7$. Chọn ngẫu nhiên ba số từ tập $E$. Xác suất gần đúng để trong ba số được chọn có đúng một số có mặt chữ số $4$ là
	 \choice
		 {$0{,}764$}
		{$0{,}125$} 
		 {$0{,}364$}
		{\True $0{,}438$} 
	\loigiai{ 
		Số các số tự nhiên có ba chữ số khác nhau từng đôi một được chọn từ các số là 
		 \[
		\mathrm{A}_8^3 - \mathrm{A}_7^2  = 294.
		\] 
		Chọn ngẫu nhiên ba số  từ tập $E \Rightarrow n(\Omega) = \mathrm{C}_{294}^3 = 4192244$.\\
		Gọi $A$ là biến cố: \lq\lq để trong ba số được chọn có đúng một số có  mặt chữ số $4$\rq\rq.\\
		Số các số tự nhiên gồm ba chữ số khác nhau từng  đôi một có mặt chữ số $4$ là
		 \[
		3 \cdot  \mathrm{A}_7^2 - 2 \cdot \mathrm{A}_6^1 = 114.
		 \]
		Số  cách chọn ba số từ tập $E$ có đúng một số có mặt chữ số $4$ là
		 \[
		\mathrm{C}_{114}^1 \cdot \mathrm{C}_{180}^2 =  1836540.
		\] 
		$$\Rightarrow \mathrm{P}(A) = \dfrac{1836540}{4192244} \approx  0.438.$$
	} 
\end{ex}

% nguồn: MaHoaNH_TracNghiem_Lop10
\begin{ex}%[0D0C2-9]%[ma:b3fb9b18]
	Một hãng hàng không phát hành $1$ triệu vé bay từ nước $X$ sang nước $Y$ có số sê-ri là một dãy gồm $6$ chữ số. Ông $A$ mua  $1$ vé bay. Sau đó, trên chuyến bay người ta thông báo vé bay may mắn là vé có số sê-ri thỏa điều kiện tổng ba chữ số đầu bằng tổng ba chữ số cuối. Nếu ai mua được vé bay may mắn sẽ được trúng thưởng một chiếc smart-phone trị giá $1\,000$ USD. Tính xác suất ông $A$ mua được vé bay may mắn. (Giả định số sê-ri được chọn hoàn toàn ngẫu nhiên).
	\choice 
		 {$0{,}015$}
		 {$0{,}225$}
		{$0{,}155$} 
		 {\True $0{,}055$}
	 \loigiai{
		Gọi dãy số may mắn có dạng $\overline{abc(9-d)(9-e)(9-f)}$ với $a,b,c,d,e,f \in  \{0,1,\dots,9\}$.\\
		 Ta có
		 \[
		a + b + c = 9 -  d + 9 - e + 9 - f \Leftrightarrow a + b + c + d + e + f = 27 \tag{*}
		\] 
		Bài toán trở thành: Đếm số nghiệm tự nhiên của phương trình $a + b + c + d + e + f = 27$ với  $0 \le a,b,c,d,e,f \le 9$.\\
		Dùng nguyên lý bù trừ (giới hạn từng biến bởi $9$), ta có số nghiệm  là
		\[ 
		 \mathrm{C}^5_{32} - \mathrm{C}^1_6 \mathrm{C}^5_{22} + \mathrm{C}^2_6 \mathrm{C}^5_{12} = 55\,252.
		\] 
		Suy ra số vé may mắn có đúng tính chất trên là $55\,252$  vé.\\
		Gọi không gian mẫu là $\Omega$, với $n(\Omega) = 10^6 = 1\,000\,000$ (tổng số sê-ri $6$ chữ số từ $000000$ đến  $999999$).\\
		Gọi $X$ là biến cố: \lq\lq Ông $A$ mua được 1 vé bay may  mắn\rq\rq, thì $n(X) = 55\,252$.\\
		 Vậy xác suất của biến cố $X$ là
		\[ 
		 \mathrm{P}(X) = \dfrac{n(X)}{n(\Omega)} = \dfrac{55\,252}{1\,000\,000} \approx 0{,}055.
		 \]
	} 
\end{ex}


%%%%%%%%%%%%%%%% TỰ LUẬN %%%%%%%%%%%%%%%%

% nguồn: MaHoaNH_TuLuan_Lop10
\begin{bt}%[0D0V2-9]%[ma:1b67b669]
	Trong trò chơi cờ tỉ phú người chơi phải gieo đồng thời hai quân súc sắc rồi sau đó cộng điểm trên các mặt của $2$ quân súc sắc để tính ra số  ô được di chuyển. Hiện tại ngay trước vị trí An đang đứng là $8$ ô tương ứng với $8$ mảnh đất của người chơi khác với giá thuê rất đắt đỏ và di chuyển thêm $10$ ô nữa thì sẽ rơi trúng ô vào tù (tính từ ô mà An đang đứng). Tính xác suất để An không phải mất tiền thuê đất và vào tù trong trò chơi này.
	\loigiai{ 
		Do gieo đồng thời hai quân súc sắc nên ta  có tổng số chấm mà An có được sẽ là\\ $\Omega=\{2;3;4;5;6;7;8;9;10;11;12\}$ nên không gian mẫu là $n(\Omega)=11$.\\
		 Gọi $A$ là biến cố: \lq \lq An không phải mất tiền thuê đất và vào tù trong trò chơi \rq \rq. \\
		Để không mất tiền thuê nhà và vào tù trong trò chơi thì tổng số chấm  trên hai quân súc sắc mà An gieo được phải lớn hơn $8$ và khác $10$.\\
		Tổng số chấm trên mặt của $2$ quân súc sắc có lợi  cho An sẽ là $9$; $11$; $12$.\\
		Khi đó số kết quả  thuận lợi cho biến cố $A$ là $n(A)=3$.\\
		Xác suất để An không mất tiền và vào tù trong trò chơi sẽ là $\mathrm{P}(A)=\dfrac{3}{11}$. 
	 }
\end{bt}

% nguồn: MaHoaNH_TuLuan_Lop10
\begin{bt}%[0D0V2-9]%[ma:2d23ecb1]
	Trong kì thi học kỳ I, bạn Bình làm để thi trắc nghiệm môn Toán. Đề thi gồm $50$ câu hỏi, mỗi câu có $4$ phương án trả lời, trong đó chỉ có một phương án đúng; trả lời đúng mỗi câu được $0{,}2$ điểm. Bình trả lời hết  các câu hỏi và chắc chắn đúng $40$ câu, $10$ câu còn lại Bình chọn ngẫu nhiên. Tính xác suất để điểm thi môn Toán của Bình không dưới $9{,}0$ điểm.
	\loigiai{Để An đúng được không  dưới $9{,}0$ điểm thì bạn ấy phải chọn đúng ít nhất $45$ câu, vậy An phải chọn đúng ít nhất $5$ trong $10$ câu còn lại.\\
		Xác suất mỗi câu trả lời đúng là $\dfrac{1}{4}$, xác suất mỗi  câu trả lời sai là $\dfrac{3}{4}$.\\
		Xác suất để An chọn đúng $5$  câu trong $10$ câu còn lại là $\mathrm{C}_{10}^{ 5} \left(\dfrac{1}{4}\right)^5\left(\dfrac{3}{4}\right)^5$.\\
		Xác suất để An chọn đúng $6$ câu trong $10$ câu còn lại là $\mathrm{C}_{10}^{ 6} \left(\dfrac{1}{4}\right)^6\left(\dfrac{3}{4}\right)^4$.\\ 
		Xác suất để An chọn đúng $7$ câu trong $10$ câu còn lại  là $\mathrm{C}_{10}^{7} \left(\dfrac{1}{4}\right)^7\left(\dfrac{3}{4}\right)^3$.\\
		Xác suất để An chọn đúng $8$ câu  trong $10$ câu còn lại là $\mathrm{C}_{10}^{8} \left(\dfrac{1}{4}\right)^8\left(\dfrac{3}{4}\right)^2$.\\
		Xác suất để An chọn đúng $9$ câu trong $10$ câu còn lại là $\mathrm{C}_{10}^{9} \left(\dfrac{1}{4}\right)^9\left(\dfrac{3}{4}\right)$.\\ 
		Xác suất để An  chọn đúng $10$ câu còn lại là $\mathrm{C}_{10}^{10} \left(\dfrac{1}{4}\right)^{10}$.\\
		Vậy  xác suất cần tìm là
		\begin{align*}
			 & \quad \mathrm{C}_{10}^{5} \left(\dfrac{1}{4}\right)^5\left(\dfrac{3}{4}\right)^5+\mathrm{C}_{10}^{6} \left(\dfrac{1}{4}\right)^6\left(\dfrac{3}{4}\right)^4+\mathrm{C}_{10}^{7} \left(\dfrac{1}{4}\right)^7\left(\dfrac{3}{4}\right)^3
			\\ &  \quad +\mathrm{C}_{10}^{8} \left(\dfrac{1}{4}\right)^8\left(\dfrac{3}{4}\right)^2+\mathrm{C}_{10}^{9} \left(\dfrac{1}{4}\right)^9\left(\dfrac{3}{4}\right)+\mathrm{C}_{10}^{10} \left(\dfrac{1}{4}\right)^{10}
			\\ &=0{,}0781. 
		\end{align*}
	} 
\end{bt}

% nguồn: MaHoaNH_TuLuan_Lop10
\begin{bt}%[0D0V2-9]%[ma:ea5ae6c2]
	Trong dịp tết Nam và Minh chơi trò ba cây bằng bộ bài lơ khơ gồm $36$ lá (mỗi người nhận được $3$ lá bài trong đó không có các quân $10$, J, Q, K sau đó cộng điểm trên $3$ quân bài lại người nào lớn điểm hơn sẽ thắng, trong trường hợp $2$ người bằng điểm thì sẽ xét chất trên các lá bài theo thứ tự từ lớn đến bé Rô, cơ, bích, tép). Người nào thắng sẽ được $3$ chiếc kẹo của người còn lại. Hiện tại trong ván bài đầu tiên Nam được $10$ điểm với các quân sau: $2$ tép, $7$ rô, Át cơ. Tính xác suất Minh  có thể thắng trong ván bài trên, chú ý Át rô là quân bài có hiệu lực mạnh nhất bộ bài (các quân Át còn lại đều tính như là số một).
	\loigiai{Do có $36$ lá bài, Nam cầm $3$ lá vì vậy còn $33$ lá bài.\\ 
		Vậy số phần tử của không gian mẫu  sẽ là $n(\Omega)=\mathrm{C}_{33}^{3}$.
		 \\
		Để thắng được Nam, Minh cần phải đạt $10$ điểm và trong đó có một quân bài chất rô lớn hơn $5$. Do Át rô  mạnh nhất nên ta xét thành $2$ trường hợp.\\
		\textbf{Trường hợp 1.} Minh có Át  rô.\\
		Như vậy để thắng được Nam, Minh cần sở  hữu $2$ lá bài còn lại sau
		$$A=\left\{ \text{$2$ bích hoặc 2 cơ hoặc 2 rô  và 7 bích hoặc 7 cơ hoặc 7 tép} \right\}.$$
		Do đó số cách ở trường hợp 1 là $3 \cdot 3=9$ cách.\\ 
		\textbf{Trường hợp 2.}Minh không có  Át rô. \\
		Như  vậy Minh bắt buộc phải có $8$ rô hoặc $9$ rô trong $3$ lá bài của mình.\\
		\textbf{Trường hợp 2.1:} Minh có  $8$ rô.\\
		Như vậy để thắng được Nam, Minh cần sở hữu $2$ lá bài còn  lại sau $$A= \left\{\text
		{$8$ bích hoặc $8$ cơ hoặc $8$ tép và $4$ bích hoặc  $4$ cơ hoặc $4$ tép hoặc $4$ rô} \right\}.$$
		Do đó số cách ở trường hợp 2.1 là  $3 \cdot 4=12$ cách.\\
		 \textbf{Trường hợp 2.2:} Minh có $9$ rô.\\
		 Như vậy để thắng được Nam, Minh cần sở hữu $2$ lá bài còn lại sau
		 $$A=\left\{\text{$9$ bích hoặc $9$ cơ hoặc $9$ tép và $2$ bích hoặc $2$ cơ hoặc $2$ rô} \right\}.$$
		Do đó số cách ở  trường hợp 2.2 là $3 \cdot 3=9$ cách.\\
		Vậy $n(A)=9+12+9=30$ nên xác suất để Minh thắng Nam sẽ  là $\mathrm{P}(A)=\dfrac{30}{5456}=\dfrac{15}{2728}$.
	} 
\end{bt}
